% =====================================================================
% New bibliography entries for the Round-1 revision (symmetry-4596052)
% Same style as the thebibliography block of manuscript.tex (MDPI/ACS):
%   Authors. Title. \textit{Journal} \textbf{Year}, \textit{Volume}, pages. https://doi.org/...
% Every DOI below was returned by Crossref (or is the arXiv-issued DOI of an
% arXiv id returned by the arXiv API). Tier C entries are commented out.
% None of these keys duplicates a key already present in manuscript.tex.
% Paste inside \begin{thebibliography}{99} ... \end{thebibliography}.
% =====================================================================

% ---------------- Tier A: must be cited -------------------------------
% Cited by the reviewers
\bibitem{ref_Xia2023} Xia, M.; B\"ottcher, L.; Chou, T. Spectrally adapted physics-informed neural networks for solving unbounded domain problems. \textit{Mach. Learn.: Sci. Technol.} \textbf{2023}, \textit{4}, 025024. https://doi.org/10.1088/2632-2153/acd0a1.
\bibitem{ref_Kutuk2025} K\"ut\"uk, M.; Y\"ucel, H. Energy dissipation preserving physics informed neural network for Allen--Cahn equations. \textit{J. Comput. Sci.} \textbf{2025}, \textit{87}, 102577. https://doi.org/10.1016/j.jocs.2025.102577.
\bibitem{ref_He2019} He, D.; Pan, K. Maximum norm error analysis of an unconditionally stable semi-implicit scheme for multi-dimensional Allen--Cahn equations. \textit{Numer. Methods Partial Differ. Equ.} \textbf{2019}, \textit{35}, 955--975. https://doi.org/10.1002/num.22333.

% Coefficient-space networks
\bibitem{ref_Du2024} Du, Y.; Chalapathi, N.; Krishnapriyan, A. Neural spectral methods: Self-supervised learning in the spectral domain. In Proceedings of the 12th International Conference on Learning Representations (ICLR), Vienna, Austria, 7--11 May 2024. https://doi.org/10.48550/arXiv.2312.05225.
\bibitem{ref_Bizzi2025} Bizzi, A.; Moreira, L.M.; Marques, M.; Mendon\c{c}a, L.; de Oliveira, C.J.; Balestro, V.; dos Santos Fernandez, L.; Yukimura, D.; Petrov, P.; Pereira, J.M.; Novello, T.; Nissenbaum, L. Neuro-spectral architectures for causal physics-informed networks. In \textit{Advances in Neural Information Processing Systems 38 (NeurIPS 2025)}; 2025; pp.\ 60600--60630. https://doi.org/10.52202/085713-1818.

% Energy-dissipative neural solvers (SAV, by construction)
\bibitem{ref_Zhang2024} Zhang, J.; Zhang, S.; Shen, J.; Lin, G. Energy-dissipative evolutionary deep operator neural networks. \textit{J. Comput. Phys.} \textbf{2024}, \textit{498}, 112638. https://doi.org/10.1016/j.jcp.2023.112638.
\bibitem{ref_Lin2025} Lin, G.; Mou, C.; Zhang, J. Energy-dissipative evolutionary Kolmogorov--Arnold networks for complex PDE systems. \textit{J. Comput. Phys.} \textbf{2025}, \textit{541}, 114326. https://doi.org/10.1016/j.jcp.2025.114326.

% Maximum bound principle
\bibitem{ref_Du2021b} Du, Q.; Ju, L.; Li, X.; Qiao, Z. Maximum bound principles for a class of semilinear parabolic equations and exponential time-differencing schemes. \textit{SIAM Rev.} \textbf{2021}, \textit{63}, 317--359. https://doi.org/10.1137/19M1243750.
\bibitem{ref_Tang2016} Tang, T.; Yang, J. Implicit-explicit scheme for the Allen--Cahn equation preserves the maximum principle. \textit{J. Comput. Math.} \textbf{2016}, \textit{34}, 451--461. https://doi.org/10.4208/jcm.1603-m2014-0017.

% PINN versus classical solvers (niche statement)
\bibitem{ref_Grossmann2024} Grossmann, T.G.; Komorowska, U.J.; Latz, J.; Sch\"onlieb, C.-B. Can physics-informed neural networks beat the finite element method? \textit{IMA J. Appl. Math.} \textbf{2024}, \textit{89}, 143--174. https://doi.org/10.1093/imamat/hxae011.

% Formal verification of numerical analysis
\bibitem{ref_Tekriwal2021} Tekriwal, M.; Duraisamy, K.; Jeannin, J.-B. A formal proof of the Lax equivalence theorem for finite difference schemes. In \textit{NASA Formal Methods}; Lecture Notes in Computer Science; Springer: Cham, Switzerland, 2021; pp.\ 322--339. https://doi.org/10.1007/978-3-030-76384-8\_20.

% ---------------- Tier B: recommended ---------------------------------
% Coefficient-space / spectral / time-marching networks
\bibitem{ref_Zhao2021} Zhao, J. Discovering phase field models from image data with the pseudo-spectral physics informed neural networks. \textit{Commun. Appl. Math. Comput.} \textbf{2021}, \textit{3}, 357--369. https://doi.org/10.1007/s42967-020-00105-2.
\bibitem{ref_Bruna2024} Bruna, J.; Peherstorfer, B.; Vanden-Eijnden, E. Neural Galerkin schemes with active learning for high-dimensional evolution equations. \textit{J. Comput. Phys.} \textbf{2024}, \textit{496}, 112588. https://doi.org/10.1016/j.jcp.2023.112588.
\bibitem{ref_Du2021a} Du, Y.; Zaki, T.A. Evolutional deep neural network. \textit{Phys. Rev. E} \textbf{2021}, \textit{104}, 045303. https://doi.org/10.1103/PhysRevE.104.045303.
\bibitem{ref_Cooley2025} Cooley, M.; Shankar, V.; Kirby, R.M.; Zhe, S. Fourier PINNs: From strong boundary conditions to adaptive Fourier bases. \textit{Trans. Mach. Learn. Res.} \textbf{2025}. https://doi.org/10.48550/arXiv.2410.03496.

% Allen--Cahn / phase-field neural solvers, energy and bounds
\bibitem{ref_Li2026} Li, C.; Yu, W.; Wang, Q. Energy dissipation rate guided adaptive sampling for physics-informed neural networks: Resolving surface-bulk dynamics in Allen--Cahn systems. \textit{Int. J. Numer. Anal. Model.} \textbf{2026}, \textit{23}, 576--604. https://doi.org/10.4208/ijnam2026-1025.
\bibitem{ref_Geng2025} Geng, Y.; Burkovska, O.; Ju, L.; Zhang, G.; Gunzburger, M. An end-to-end deep learning method for solving nonlocal Allen--Cahn and Cahn--Hilliard phase-field models. \textit{Comput. Methods Appl. Mech. Eng.} \textbf{2025}, \textit{436}, 117721. https://doi.org/10.1016/j.cma.2024.117721.
\bibitem{ref_Shen2016} Shen, J.; Tang, T.; Yang, J. On the maximum principle preserving schemes for the generalized Allen--Cahn equation. \textit{Commun. Math. Sci.} \textbf{2016}, \textit{14}, 1517--1534. https://doi.org/10.4310/CMS.2016.v14.n6.a3.

% Exact (hard) constraints and structure by construction
\bibitem{ref_Dong2021} Dong, S.; Ni, N. A method for representing periodic functions and enforcing exactly periodic boundary conditions with deep neural networks. \textit{J. Comput. Phys.} \textbf{2021}, \textit{435}, 110242. https://doi.org/10.1016/j.jcp.2021.110242.
\bibitem{ref_Sukumar2022} Sukumar, N.; Srivastava, A. Exact imposition of boundary conditions with distance functions in physics-informed deep neural networks. \textit{Comput. Methods Appl. Mech. Eng.} \textbf{2022}, \textit{389}, 114333. https://doi.org/10.1016/j.cma.2021.114333.
\bibitem{ref_Lu2021} Lu, L.; Pestourie, R.; Yao, W.; Wang, Z.; Verdugo, F.; Johnson, S.G. Physics-informed neural networks with hard constraints for inverse design. \textit{SIAM J. Sci. Comput.} \textbf{2021}, \textit{43}, B1105--B1132. https://doi.org/10.1137/21M1397908.
\bibitem{ref_CardosoBihlo2025} Cardoso-Bihlo, E.; Bihlo, A. Exactly conservative physics-informed neural networks and deep operator networks for dynamical systems. \textit{Neural Netw.} \textbf{2025}, \textit{181}, 106826. https://doi.org/10.1016/j.neunet.2024.106826.

% Failure modes, conditioning, baselines, error estimates
\bibitem{ref_Krishnapriyan2021} Krishnapriyan, A.S.; Gholami, A.; Zhe, S.; Kirby, R.M.; Mahoney, M.W. Characterizing possible failure modes in physics-informed neural networks. In \textit{Advances in Neural Information Processing Systems 34 (NeurIPS 2021)}; 2021. https://doi.org/10.48550/arXiv.2109.01050.
\bibitem{ref_Rathore2024} Rathore, P.; Lei, W.; Frangella, Z.; Lu, L.; Udell, M. Challenges in training PINNs: A loss landscape perspective. In Proceedings of the 41st International Conference on Machine Learning (ICML); \textit{Proc. Mach. Learn. Res.} \textbf{2024}, \textit{235}, pp.\ 42159--42191. https://doi.org/10.48550/arXiv.2402.01868.
\bibitem{ref_Wang2021a} Wang, S.; Teng, Y.; Perdikaris, P. Understanding and mitigating gradient flow pathologies in physics-informed neural networks. \textit{SIAM J. Sci. Comput.} \textbf{2021}, \textit{43}, A3055--A3081. https://doi.org/10.1137/20M1318043.
\bibitem{ref_Rahaman2019} Rahaman, N.; Baratin, A.; Arpit, D.; Draxler, F.; Lin, M.; Hamprecht, F.A.; Bengio, Y.; Courville, A. On the spectral bias of neural networks. In Proceedings of the 36th International Conference on Machine Learning (ICML); \textit{Proc. Mach. Learn. Res.} \textbf{2019}, \textit{97}, pp.\ 5301--5310. https://doi.org/10.48550/arXiv.1806.08734.
\bibitem{ref_McGreivy2024} McGreivy, N.; Hakim, A. Weak baselines and reporting biases lead to overoptimism in machine learning for fluid-related partial differential equations. \textit{Nat. Mach. Intell.} \textbf{2024}, \textit{6}, 1256--1269. https://doi.org/10.1038/s42256-024-00897-5.
\bibitem{ref_Mishra2023} Mishra, S.; Molinaro, R. Estimates on the generalization error of physics-informed neural networks for approximating PDEs. \textit{IMA J. Numer. Anal.} \textbf{2023}, \textit{43}, 1--43. https://doi.org/10.1093/imanum/drab093.

% Formal verification of numerical analysis
\bibitem{ref_Boldo2013} Boldo, S.; Cl\'ement, F.; Filli\^atre, J.-C.; Mayero, M.; Melquiond, G.; Weis, P. Wave equation numerical resolution: A comprehensive mechanized proof of a C program. \textit{J. Autom. Reason.} \textbf{2013}, \textit{50}, 423--456. https://doi.org/10.1007/s10817-012-9255-4.
\bibitem{ref_Boldo2017} Boldo, S.; Cl\'ement, F.; Faissole, F.; Martin, V.; Mayero, M. A Coq formal proof of the Lax--Milgram theorem. In Proceedings of the 6th ACM SIGPLAN Conference on Certified Programs and Proofs (CPP); ACM: New York, NY, USA, 2017; pp.\ 79--89. https://doi.org/10.1145/3018610.3018625.
\bibitem{ref_Immler2018} Immler, F. A verified ODE solver and the Lorenz attractor. \textit{J. Autom. Reason.} \textbf{2018}, \textit{61}, 73--111. https://doi.org/10.1007/s10817-017-9448-y.

% ---------------- Tier C: optional (arXiv preprints, not known to be peer reviewed) -------------
% Uncomment only if you want them; cite as preprints.
% \bibitem{ref_Marsh2026} Marsh, C.; Erban, R.; Munch, A. Extended pseudo-spectral physics-informed neural networks for phase-field models. \textit{arXiv} \textbf{2026}, arXiv:2606.24660. https://doi.org/10.48550/arXiv.2606.24660.
% \bibitem{ref_Sivalingam2025} Sivalingam, S.M.; Govindaraj, V.; Hendy, A.S. Spectral coefficient learning physics informed neural network for time-dependent fractional parametric differential problems. \textit{arXiv} \textbf{2025}, arXiv:2503.22386. https://doi.org/10.48550/arXiv.2503.22386.
% \bibitem{ref_Nguyen2026} Nguyen, T.-P.; Tran, M.-B.; Tu, S. A structure-preserving neural-spectral method for reconstructing controls of wave equations. \textit{arXiv} \textbf{2026}, arXiv:2606.28029. https://doi.org/10.48550/arXiv.2606.28029.
% \bibitem{ref_Yin2024} Yin, P.; Ling, S.; Ying, W. Chebyshev spectral neural networks for solving partial differential equations. \textit{arXiv} \textbf{2024}, arXiv:2407.03347. https://doi.org/10.48550/arXiv.2407.03347.
% \bibitem{ref_Tang2026} Tang, T.; Yang, J.; Zhao, Y.; Zhu, Q. Energy dissipation preserving feature-based DNN Galerkin methods for gradient flows. \textit{arXiv} \textbf{2026}, arXiv:2603.14029. https://doi.org/10.48550/arXiv.2603.14029.
% \bibitem{ref_Zhang2026} Zhang, Z.; Alla, A.; Karniadakis, G.E. Adjoint method versus physics-informed neural networks in PDE-constrained inverse problems. \textit{arXiv} \textbf{2026}, arXiv:2606.12337. https://doi.org/10.48550/arXiv.2606.12337.
% \bibitem{ref_Wang2021b} Wang, S.; Wang, H.; Perdikaris, P. On the eigenvector bias of Fourier feature networks: From regression to solving multi-scale PDEs with physics-informed neural networks. \textit{Comput. Methods Appl. Mech. Eng.} \textbf{2021}, \textit{384}, 113938. https://doi.org/10.1016/j.cma.2021.113938.

% ---------------------------------------------------------------------
% Existing keys in manuscript.tex that are the SAME works as items above: none.
% (The three reviewer-cited works Xia2023, Kutuk2025 and He2019 are not yet in the bibliography.)
% Existing keys to cite alongside the new ones (do not duplicate):
%   Guo_EnergySeq2024, Wight, Mattey, Mostajeran_cPIKAN2025, Zhu_SPINN2022, Zhang_SPINN2023,
%   Zhang_SymGroup2023, Lu_DeepONet, Li_FNO, Wang_NTK2022, Wang_causal, Liu_FourierBias2025,
%   Shen_SAV, Zhang_ETD2024, Gao_pIFRK2023, Du_MPPDG2022, Feng_MPPFD2021, Kim_MPP3D2025,
%   DeMouraUllrich2021, Bobbin2023, Ugwuanyi2025, Tsachouridis2020.
% ---------------------------------------------------------------------
